% !TeX spellcheck = en_US
%!TEX root = ../Main.tex

\section{Introduction}
\label{sec:introduction}

Foundation models have been gaining increasing attention in the research community \cite{bommasani2021opportunities,akbari2021vatt,shao2021intern}, since they give a practical paradigm for scaling to numerous perception tasks with surprisingly good results. Through simple adaption or zero/few-shot learning, foundation models greatly reduce downstream design and training costs using generic representations learned from web-scale data with a strong backbone of high capacity. It is expected that developing foundation models can cultivate cognition from perception, obtaining general vision capability.

Though a line of vision foundation models is proposed \cite{CLIP,align,florence,unicl,simvlm,ofa,coca,beit3,beit,barham2022pathways}, video understanding and the corresponding tasks are less explored compared with image ones, mainly used for validating that the visual features from these models are also beneficial to spatiotemporal representations. We conjecture this relatively scarce focus from the academic community is caused by 1) a high computing burden from video processing, and 2) quite a few current video benchmarks can be handled by exploiting appearance features from image backbones with accordingly temporal modeling. Specifically, for efficiency, the additional time dimension in video processing makes at least one order of magnitude higher complexity than image processing when their spatial resolutions are close and the temporal sampling ratio is usually 16. 
For some current video datasets, image features alone or with lateral temporal modules are sufficient to give decent results, especially with the rise of the multimodal model CLIP \cite{CLIP}. Its various temporal variants yield competitive or state-of-the-art performance in several core tasks \cite{clip4clip,ma2022x}. 
%We attribute this to the missing of a proper comprehensive video understanding benchmark, especially temporal understanding ones. 
Regarding this, a simultaneous spatiotemporal learner does not seem like a sweet spot between research \& development cost and payback.

Moreover, the transferability of current vision foundation models is somewhat narrow considering the wide spectrum of video applications. These models~\cite{videomae,mtv,maskfeat,all_in_one} either concentrate on action understanding tasks (\eg action recognition, spatiotemporal action localization, \etc) or video-language alignment ones (\eg video retrieval, video question answering, \etc). We suppose this results from their learning schemes, as well as the lack of a comprehensive benchmark for measuring video understanding capabilities. Thus, these works~\cite{videomae,mtv,maskfeat,all_in_one} focalize a few specific tasks to demonstrate their spatiotemporal perceptions.
The community desires a general foundation model that enables a broader application domain.

In this paper, we advance video foundation model research with a cost-effective and versatile model \model. 
To establish a feasible and effective spatiotemporal representation, we study both popular video masked modeling \cite{mae,videomae} and multimodal contrastive learning \cite{CLIP,cpd}. Note that video masking modeling specializes in action understanding, and it is still worth exploring regarding its limited model scale caused by the current decoder. For multimodal contrastive learning, it embeds rich semantics into video representation while ignoring concrete spatiotemporal modeling. To address these challenges, we make these two self-supervised approaches learn at scale efficiently in modular designs.
To significantly broaden the generalization of the current video foundation models, we propose a unified representation learning with both two self-supervised training manners. To validate such a generalized representation, we propose a systematic video understanding benchmark. It involves evaluations of action understanding, video-language alignment, and open-world video applications, which we believe are three core abilities of generic video perception. Instead of introducing new data or annotations to this system, we initially choose ten representative video tasks with 39 public datasets, and categorize them into those three types. 
To our best knowledge, \model~is the first video foundation model which demonstrates promising transferability with state-of-the-art performance in all those three different types of video tasks.

In \model, we design a unified video representation (UVR) learning paradigm. It explores both masked video modeling with autoencoders (MAE) and multimodal contrastive learning for two types of representations, strengthens them by supervised action classification, and generates a more general representation based on the cross-representation learning between them.
UVR not only empirically shows video representation outperforms image one with temporal capturing significantly on core video tasks, but also is training-efficient. Its MAE exploits high redundancies in videos and trains with only a few visible tokens. Meanwhile, multimodal learning in \model~extends existing image-pretrained backbones for video contrastive training. After supervised training these two video encoders, we craft cross-model attention to conduct feature alignments between these two almost frozen encoders.

More than a unified video representation learning paradigm, we also make practices and guidelines for training large-scale video foundation models in a tractable and efficient manner. Our work contains and is not limited to 1) making VideoMAE scalable and exploring its scalability in model and data scale; 2) efficient and effective multimodal architecture design and training receipt about how to leverage existing image-pretrained backbones; 3) empirically finding features from VideoMAE and multimodal models are complementary and studying how to deduce more powerful video representations by coordinating different existing models. Specifically, 

\begin{itemize}[leftmargin=*]
	
	\item For the scalability study of VideoMAE, we show that the proper diversity and scaling-up size in training videos can improve the scalability of the used video encoder. With a new pretrained dataset in a masked autoencoder training setting, ViT lifts its action recognition performance on Kinetics-400 \cite{k400} with finetuning from 81.01\% to 85.35\% from base to large, and further reaches 86.9\% with the huge setup, surpassing the performance reported in \cite{videomae} with a notable margin. The scalability of VideoMAE enables its usage in video foundation model development.
	
	\item For reusing existing foundation models for multimodal learning, we extend an image-pretrained vision transformer \cite{dosovitskiy2020image} to video representation learning. This transfer learning requires substantial structure and optimization customizations or using local and global spatiotemporal modules for multimodal pretraining. The local module disentangles spatiotemporal modeling by consecutive and independent spatial and temporal attention computation. Meanwhile, the global module computes token interaction across space and time. Experiments show that this reuse design is effective in spatiotemporal representation learning.
	
	\item More than self-supervised pretraining, we also employ supervised action recognition to further enhance video representation. Results demonstrate action recognition is a fine source task for transferring to various downstream applications.
	
	\item To coordinate foundation models, we unify masked video encoder with multimodal one by cross-representation learning, instead of training them jointly in one formulation. Regarding the optimization of MAE and multimodal learning (MML) that may contradict each other, combining them without compromising their merits remains an open question \cite{bachmann2022multimae}. More importantly, MML with contrastive learning demands huge batches for better contrastive optimization. Adding MAE to it will inevitably lead to numerous headachy implementation issues. Considering their potential training adversaries, we train MAE and MML separately. After their training converges, we then dynamically combine their representations with the proposed cross-model attention (CMA) modules. It implements cross-attention between MAE and MML mid-level features, adaptively fusing their high-level features for prediction. In the model-level representation interaction phase, we freeze backbones trained by MAE and MML separately and only let CMA be updated in supervised learning with a few epochs. Experiments show it is a computationally tractable and efficient means to exploit both MAE and MML features. 
\end{itemize}

We validate our proposed video foundation model in 10 tasks with 39 datasets (including core tasks \eg\ action recognition, spatiotemporal action localization, video question answering, video retrieval, \etc), and it outperforms all the state-of-the-art methods in each task non-trivially. We suppose these overall superior results obtained by our approach, along with observations and analysis, set up a new baseline for the video understanding community. The empirical evidence in this paper raises the confidence that video perceptive tasks and partial high-order tasks (formulated to perceptive forms) can be well-addressed by video foundation models, serving as a performance-critical method across a spectrum of applications.

In summary, we contribute to video foundation models in the following aspects:
\begin{itemize}[leftmargin=*]
	\item We explore a general video representation paradigm with both masked and contrastive modeling, and realize this design by unifying their representation by lightweight model interaction learning in supervision. We confirm features learned by generative and contrastive training are complementary to experiments and can deliver better results than either of them trained independently.
	\item We find masked video encoder can be scalable in model and data size with proper tuning. We devise pluggable local temporal and global spatiotemporal interaction modules to reuse pretrained ViT with image-text data for multimodal learning, easing the training burden and yielding better downstream performance.
	\item We make a tentative attempt in constructing a systematic video understanding benchmark. Our general video foundation models achieve state-of-the-art performance on 39 datasets with several core tasks in this benchmark, \eg, Kinetics-400 and Something-Something v2 in action recognition. We empirically find our learned video representations outperform their rivals, dominating vision-language tasks by a large margin, especially for some image-based ones. It suggests general video representations will be a central player in video tasks. We believe the openness of our proposed methods and models will provide the research community with handy tools to foundation models and their features with easy access. 
\end{itemize}
